\documentclass[12pt,a4paper]{article}
\usepackage[a4paper,margin=2cm]{geometry}
\usepackage[T2A]{fontenc}
\usepackage[utf8]{inputenc}
\usepackage[russian]{babel}
\usepackage{amsmath,amssymb,mathtools}
\usepackage{array,booktabs,longtable}
\usepackage{xcolor}
\usepackage{hyperref}
\usepackage{tikz}
\setlength{\parindent}{0pt}
\setlength{\parskip}{0.55em}
\hypersetup{colorlinks=true,linkcolor=blue!55!black,urlcolor=blue!55!black}
\definecolor{errred}{RGB}{175,35,35}
\definecolor{ideablue}{RGB}{35,75,135}
\definecolor{resultteal}{RGB}{20,105,95}

\makeatletter
\def\ps@reviewstyle{%
  \def\@oddhead{}\def\@evenhead{}%
  \def\@oddfoot{\hfil\small Лёвин Артём Александрович эксклюзивно для Анастасии Павловой\hfil}%
  \let\@evenfoot\@oddfoot}
\makeatother
\pagestyle{reviewstyle}

\newcommand{\errorlabel}[1]{\textcolor{errred}{\textbf{Ошибка:}} \textbf{#1}}
\newcommand{\keyidea}[1]{\textcolor{ideablue}{\textbf{Ключевая идея:}} #1}
\newcommand{\finalresult}[1]{\textcolor{resultteal}{\textbf{Итог:}} #1}

\begin{document}

% Краткое сообщение Анастасии: Анастасия, в работе есть ошибки в переводе единиц, подстановке данных из условия, арифметическом вычислении, согласовании минут и часов и раскрытии скобок; последнее задание не доведено до итогового ответа.


\begin{titlepage}
\thispagestyle{reviewstyle}
\centering
\vspace*{0.22\textheight}
{\Huge\bfseries Ревью домашней работы\par}
\vspace{1.8cm}
{\Large Автор: Лёвин Артём Александрович\par}
\vspace{0.8cm}
{\large 5 октября 2026 года\par}
\vfill
\begin{tikzpicture}[scale=1]
\draw[line width=0.8pt,blue!35!black] (0,0) .. controls (2,0.7) and (4,-0.7) .. (6,0);
\end{tikzpicture}
\end{titlepage}

\newpage
\section*{Сводная таблица}
\small
\renewcommand{\arraystretch}{1.25}
\setlength{\tabcolsep}{3pt}
\begin{longtable}{>{\raggedright\arraybackslash}p{0.08\linewidth} >{\raggedright\arraybackslash}p{0.22\linewidth} >{\raggedright\arraybackslash}p{0.22\linewidth} >{\raggedright\arraybackslash}p{0.16\linewidth} >{\raggedright\arraybackslash}p{0.17\linewidth}}
\toprule
№ задания & Тема & Суть ошибки & Ваш ответ & Правильный ответ \\
\midrule
\endfirsthead
\toprule
№ задания & Тема & Суть ошибки & Ваш ответ & Правильный ответ \\
\midrule
\endhead
\hyperref[task:1]{1} & Движение мимо неподвижного объекта & Неверно переведены единицы времени и скорости; результат завышен в десять раз. & 5000 м & 500 м \\
\hyperref[task:3]{3} & Встречное движение поездов & В расчёт подставлены данные не из условия задачи: другие длины и скорости. & 7,5 с & $\dfrac{80}{7}$ с $\approx 11{,}43$ с \\
\hyperref[task:5]{5} & Относительное движение & Правильно составлено выражение для относительного пути, но допущена десятикратная арифметическая ошибка при делении. & 30 км/ч & 3 км/ч \\
\hyperref[task:8]{8} & Совместная величина и разность времени & В одном уравнении смешаны часы и минуты: разность времени взята как 18, хотя скорости заданы в вопросах в час. & 0,225 & 12 вопросов \\
\hyperref[task:9]{9} & Производительность труб & При раскрытии скобок потерян множитель 600 перед переменной. & около 5,26 л/мин & 40 л/мин \\
\hyperref[task:10]{10} & Совместная работа & Решение не завершено: после записи исходных времён не найден объём выполненной и оставшейся работы. & не указан & 2 ч 48 мин \\
\bottomrule
\end{longtable}
\normalsize

\newpage
\vspace*{0.35\textheight}
\begin{center}
{\LARGE\bfseries Разбор ошибок}
\end{center}
\newpage

\phantomsection\label{task:1}
\section*{Задание №1}
\textbf{Текст задания.} Поезд движется со скоростью 72 км/ч и проходит мимо столба за 25 секунд. Найти длину поезда в метрах.

\textbf{Ваше решение.} В записи использована формула движения и затем получено значение 5 км, после чего оно переведено как 5000 м.

\textbf{Ваш ответ.} 5000 м.

\errorlabel{неверное согласование единиц измерения.}

\textbf{Где именно возникает ошибка.} Саму идею использовать связь между расстоянием, скоростью и временем можно сохранить. Ошибка появляется в момент, когда скорость 72 км/ч и время 25 секунд используются без корректного приведения к одной системе единиц. Из-за этого длина получается в десять раз больше правильной.

\textbf{Почему это неверно.} Если скорость дана в километрах в час, а время --- в секундах, нельзя непосредственно перемножать эти величины. Нужно либо перевести скорость в метры в секунду, либо перевести 25 секунд в часы. Самый удобный школьный путь здесь --- перевести скорость: один километр содержит 1000 метров, один час содержит 3600 секунд, поэтому
\[
72\ \text{км/ч}=72\cdot\frac{1000}{3600}\ \text{м/с}=20\ \text{м/с}.
\]
После этого скорость и время согласованы: метры в секунду умножаются на секунды, и получается расстояние в метрах.

\keyidea{Когда поезд проходит мимо столба, за время прохождения он перемещается ровно на собственную длину. Поэтому длина поезда равна произведению скорости на время, но только после приведения единиц к одной системе.}

\textbf{Исправление от последнего правильного шага.} Используем формулу длины через скорость и время:
\[
L=vt.
\]
Переводим скорость:
\[
72\ \text{км/ч}=20\ \text{м/с}.
\]
Тогда
\[
L=20\cdot25=500\ \text{м}.
\]

\textbf{Полное правильное решение.} За одну секунду поезд проходит 20 метров. За 25 секунд он проходит
\[
20\cdot25=500\ \text{м}.
\]
Поскольку столб является неподвижной точкой, для полного прохождения мимо него поезд должен сместиться на расстояние, равное своей длине. Следовательно, длина поезда равна 500 м.

\textbf{Проверка.} При скорости 20 м/с за 10 секунд поезд проходит 200 м, за 20 секунд --- 400 м, ещё за 5 секунд --- 100 м. Итого 500 м, что согласуется с вычислением.

\finalresult{длина поезда равна 500 м.}

\subsection*{Задача на закрепление}
Поезд движется со скоростью 90 км/ч и проходит мимо столба за 18 секунд. Найдите длину поезда в метрах.

\newpage
\phantomsection\label{task:3}
\section*{Задание №3}
\textbf{Текст задания.} Два поезда длиной 180 м и 220 м движутся навстречу друг другу со скоростями 54 км/ч и 72 км/ч. Сколько секунд проходит от встречи их локомотивов до полного расхождения составов?

\textbf{Ваше решение.} В записи использовано сложение длин и скоростей, однако подставлены величины 0,12 км и 0,18 км, а также скорости 90 км/ч и 54 км/ч. Получен ответ 7,5 с.

\textbf{Ваш ответ.} 7,5 с.

\errorlabel{в расчёт перенесены данные из другого условия.}

\textbf{Где именно возникает ошибка.} Идея решения верная: при встречном движении нужно сложить длины поездов и сложить их скорости. Первый неверный шаг --- подстановка в эту схему длин 120 м и 180 м и скоростей 90 км/ч и 54 км/ч вместо данных задачи №3.

\textbf{Почему это неверно.} От встречи локомотивов до полного расхождения каждый состав должен пройти относительно другого суммарную длину двух поездов. В этой задаче суммарная длина равна 180 м плюс 220 м, то есть 400 м. Относительная скорость при движении навстречу равна сумме скоростей 54 км/ч и 72 км/ч. Если заменить хотя бы одну из этих величин данными другой задачи, вычисляется уже другая ситуация и итоговое время теряет связь с исходным условием.

\keyidea{Для двух объектов, движущихся навстречу, относительная скорость равна сумме их скоростей. Для полного расхождения поездов относительно друг друга нужно пройти сумму их длин.}

\textbf{Исправление от последнего правильного шага.} Сохраняем Вашу схему, но подставляем правильные данные:
\[
L_{\text{общ}}=180+220=400\ \text{м}.
\]
Переведём скорости в метры в секунду:
\[
54\ \text{км/ч}=15\ \text{м/с},\qquad 72\ \text{км/ч}=20\ \text{м/с}.
\]
Относительная скорость:
\[
v_{\text{отн}}=15+20=35\ \text{м/с}.
\]
Тогда
\[
t=\frac{400}{35}=\frac{80}{7}\ \text{с}\approx 11{,}43\ \text{с}.
\]

\textbf{Полное правильное решение.} После встречи локомотивов поезда начинают проходить один относительно другого. Полное расхождение произойдёт, когда относительное перемещение станет равно сумме длин обоих составов:
\[
180+220=400\ \text{м}.
\]
Скорости переводим в метры в секунду:
\[
54:3{,}6=15\ \text{м/с},\qquad 72:3{,}6=20\ \text{м/с}.
\]
Так как поезда движутся навстречу, их относительная скорость равна
\[
15+20=35\ \text{м/с}.
\]
По формуле времени:
\[
t=\frac{400}{35}=\frac{80}{7}\ \text{с}.
\]

\textbf{Проверка.} За примерно 11,43 секунды при относительной скорости 35 м/с поезда проходят относительно друг друга примерно $35\cdot11{,}43\approx400$ м, то есть ровно суммарную длину составов.

\finalresult{от встречи локомотивов до полного расхождения проходит $\dfrac{80}{7}$ с, то есть примерно 11,43 с.}

\subsection*{Задача на закрепление}
Два поезда длиной 150 м и 200 м движутся навстречу друг другу со скоростями 72 км/ч и 54 км/ч. Сколько секунд проходит от встречи локомотивов до полного расхождения составов?

\newpage
\phantomsection\label{task:5}
\section*{Задание №5}
\textbf{Текст задания.} Две баржи длиной 70 м и 50 м движутся в одном направлении. До начала контакта более быстрой барже нужно сократить расстояние 250 м, а после полного обгона пройти ещё 180 м относительно медленной. На весь манёвр ушло 11 минут. На сколько км/ч скорость быстрой баржи больше скорости медленной?

\textbf{Ваше решение.} Правильно учтены четыре части относительного пути: 70 м, 50 м, 250 м и 180 м. Получена сумма 0,55 км и использовано время 11 минут, то есть $11/60$ часа. Затем получен ответ 30 км/ч.

\textbf{Ваш ответ.} 30 км/ч.

\errorlabel{десятикратная арифметическая ошибка в последнем вычислении.}

\textbf{Где именно возникает ошибка.} До последнего арифметического действия решение построено правильно. Из равенства
\[
\frac{11}{60}=\frac{0{,}55}{v_{\text{быстр}}-v_{\text{медл}}}
\]
следует
\[
v_{\text{быстр}}-v_{\text{медл}}=\frac{0{,}55\cdot60}{11}.
\]
Именно при вычислении этого выражения появляется неверное значение 30 вместо 3.

\textbf{Почему это неверно.} Произведение $0{,}55\cdot60$ равно 33. После деления 33 на 11 получается 3. Десятичная запятая здесь принципиальна: число 0,55 меньше единицы, поэтому результат не может внезапно стать в десять раз больше без дополнительного множителя.

\keyidea{При обгоне в одном направлении удобно перейти к относительному движению. Медленная баржа как будто неподвижна, а быстрая движется относительно неё со скоростью, равной разности скоростей.}

\textbf{Исправление от последнего правильного шага.} Относительный путь равен
\[
70+50+250+180=550\ \text{м}=0{,}55\ \text{км}.
\]
Время:
\[
11\ \text{мин}=\frac{11}{60}\ \text{ч}.
\]
Тогда разность скоростей:
\[
\Delta v=\frac{0{,}55}{11/60}=0{,}55\cdot\frac{60}{11}=\frac{33}{11}=3\ \text{км/ч}.
\]

\textbf{Полное правильное решение.} До соприкосновения быстрая баржа должна сократить 250 м. Чтобы полностью пройти вдоль медленной баржи, нужно относительно неё пройти сумму длин обеих барж:
\[
70+50=120\ \text{м}.
\]
После полного обгона по условию она проходит ещё 180 м относительно медленной. Поэтому общий относительный путь:
\[
250+120+180=550\ \text{м}=0{,}55\ \text{км}.
\]
Время манёвра равно 11 минутам, или $11/60$ часа. Разность скоростей равна отношению относительного пути ко времени:
\[
\Delta v=\frac{0{,}55}{11/60}=3\ \text{км/ч}.
\]

\textbf{Проверка.} При разности скоростей 3 км/ч за $11/60$ часа относительный путь равен
\[
3\cdot\frac{11}{60}=\frac{33}{60}=0{,}55\ \text{км},
\]
то есть 550 м. Это точно совпадает с требуемым относительным путём.

\finalresult{быстрая баржа движется на 3 км/ч быстрее медленной.}

\subsection*{Задача на закрепление}
Две баржи длиной 80 м и 40 м движутся в одном направлении. До начала контакта более быстрой барже нужно сократить расстояние 300 м, а после полного обгона пройти ещё 180 м относительно медленной. На весь манёвр ушло 12 минут. На сколько км/ч скорость быстрой баржи больше скорости медленной?

\newpage
\phantomsection\label{task:8}
\section*{Задание №8}
\textbf{Текст задания.} Петя отвечает на 8 вопросов теста в час, Ваня --- на 10 вопросов в час. Тест у них одинаковый. Ваня заканчивает на 18 минут раньше Пети. Сколько вопросов в тесте?

\textbf{Ваше решение.} Обозначено количество вопросов через $x$. Время Пети записано как $x/8$, время Вани --- как $x/10$. Затем составлено равенство разности этих времён числу 18.

\textbf{Ваш ответ.} 0,225.

\errorlabel{в одном уравнении смешаны часы и минуты.}

\textbf{Где именно возникает ошибка.} Записи $x/8$ и $x/10$ верны: поскольку производительности заданы в вопросах за час, обе эти величины измеряются в часах. Первый неверный шаг --- приравнивание их разности к 18. Число 18 в условии означает 18 минут, а не 18 часов.

\textbf{Почему это неверно.} Уравнение должно сравнивать величины в одинаковых единицах. Разность $x/8-x/10$ измеряется в часах. Поэтому 18 минут нужно перевести в часы:
\[
18\ \text{мин}=\frac{18}{60}\ \text{ч}=0{,}3\ \text{ч}.
\]
Только после этого можно составлять уравнение.

\keyidea{Когда производительность дана «в час», все времена в уравнении удобно выражать в часах. Перевод минут нужно выполнить до составления уравнения.}

\textbf{Исправление от последнего правильного шага.} Сохраняем Ваши выражения времени:
\[
t_{\text{Пети}}=\frac{x}{8},\qquad t_{\text{Вани}}=\frac{x}{10}.
\]
Так как Ваня заканчивает на 18 минут, то есть на 0,3 часа раньше,
\[
\frac{x}{8}-\frac{x}{10}=0{,}3.
\]
Приводим левую часть к общему знаменателю:
\[
\frac{5x-4x}{40}=0{,}3,
\]
\[
\frac{x}{40}=0{,}3,
\]
\[
x=12.
\]

\textbf{Полное правильное решение.} Пусть в тесте $x$ вопросов. Петя отвечает на 8 вопросов за час, значит ему нужно $x/8$ часа. Ваня отвечает на 10 вопросов за час, значит ему нужно $x/10$ часа. Разность времени равна 18 минутам, то есть 0,3 часа:
\[
\frac{x}{8}-\frac{x}{10}=0{,}3.
\]
Разность дробей равна $x/40$, поэтому
\[
\frac{x}{40}=0{,}3.
\]
Отсюда
\[
x=12.
\]

\textbf{Проверка.} Петя решит 12 вопросов за $12/8=1{,}5$ часа, то есть за 90 минут. Ваня решит 12 вопросов за $12/10=1{,}2$ часа, то есть за 72 минуты. Разность равна 18 минутам, как в условии.

\finalresult{в тесте 12 вопросов.}

\subsection*{Задача на закрепление}
Один участник отвечает на 12 вопросов в час, другой --- на 15 вопросов в час. Тест у них одинаковый, и второй заканчивает на 20 минут раньше первого. Сколько вопросов в тесте?

\newpage
\phantomsection\label{task:9}
\section*{Задание №9}
\textbf{Текст задания.} Две трубы работают одинаковое время. Первая наполняет бак объёмом 600 л, вторая --- бак объёмом 690 л. Производительность второй трубы на 6 л/мин больше производительности первой. Найти производительность первой трубы.

\textbf{Ваше решение.} Производительность первой трубы обозначена через $x$, второй --- через $x+6$. Составлено верное равенство времени
\[
\frac{600}{x}=\frac{690}{x+6}.
\]
После умножения на знаменатели записан переход, в котором выражение $600(x+6)$ фактически превращено в $x+3600$.

\textbf{Ваш ответ.} В записи получено значение около $\dfrac{720}{137}$ л/мин.

\errorlabel{потерян множитель при раскрытии скобок.}

\textbf{Где именно возникает ошибка.} Последний полностью правильный шаг:
\[
\frac{600}{x}=\frac{690}{x+6}.
\]
После перекрёстного умножения должно получиться
\[
600(x+6)=690x.
\]
Первый неверный переход --- раскрытие левой части как $x+3600$ вместо $600x+3600$.

\textbf{Почему это неверно.} По распределительному закону множитель перед скобками умножается на каждый член внутри скобок:
\[
600(x+6)=600\cdot x+600\cdot6=600x+3600.
\]
Нельзя умножить 600 только на число 6 и не умножить на $x$. Потеря множителя резко меняет уравнение и приводит к неверному значению производительности.

\keyidea{В задачах на производительность время равно объёму, делённому на производительность. Если времена равны, соответствующие дроби можно приравнять, а при раскрытии скобок нужно распределить внешний множитель на каждый член.}

\textbf{Исправление от последнего правильного шага.}
\[
\frac{600}{x}=\frac{690}{x+6}.
\]
Перемножаем крест-накрест:
\[
600(x+6)=690x.
\]
Раскрываем скобки правильно:
\[
600x+3600=690x.
\]
Переносим $600x$ вправо:
\[
3600=90x.
\]
Отсюда
\[
x=40.
\]

\textbf{Полное правильное решение.} Пусть производительность первой трубы равна $x$ л/мин. Тогда производительность второй трубы равна $x+6$ л/мин. Первая труба наполняет 600 л за $600/x$ минут, вторая --- 690 л за $690/(x+6)$ минут. По условию времена одинаковы:
\[
\frac{600}{x}=\frac{690}{x+6}.
\]
После перекрёстного умножения:
\[
600x+3600=690x.
\]
Следовательно,
\[
3600=90x,
\]
\[
x=40.
\]
Значит, первая труба подаёт 40 л в минуту, а вторая --- 46 л в минуту.

\textbf{Проверка.} Первая труба заполнит 600 л за
\[
600:40=15\ \text{мин}.
\]
Вторая труба заполнит 690 л за
\[
690:46=15\ \text{мин}.
\]
Времена действительно одинаковы. Разность производительностей равна 6 л/мин, как в условии.

\finalresult{производительность первой трубы равна 40 л/мин.}

\subsection*{Задача на закрепление}
Две трубы работают одинаковое время. Первая наполняет 480 л, вторая --- 600 л. Производительность второй трубы на 8 л/мин больше производительности первой. Найдите производительность первой трубы.

\newpage
\phantomsection\label{task:10}
\section*{Задание №10}
\textbf{Текст задания.} Первый мастер выполняет заказ за 5 часов, второй --- за 8 часов. Они работают вместе 2 часа, после чего первый мастер уходит. Сколько времени потребуется второму мастеру, чтобы закончить оставшуюся работу?

\textbf{Ваше решение.} Начата таблица со временем выполнения заказа: 5 часов для первого мастера и 8 часов для второго. Дальнейшие вычисления и итоговый ответ отсутствуют.

\textbf{Ваш ответ.} Не указан.

\errorlabel{решение не доведено до вычисления оставшейся работы и времени.}

\textbf{Где именно возникает ошибка.} Записанные исходные времена 5 часов и 8 часов верны. Следующий обязательный шаг --- перейти от времени полного заказа к производительности каждого мастера: первый выполняет одну пятую заказа за час, второй --- одну восьмую заказа за час. Этот шаг и последующее вычисление не выполнены.

\textbf{Почему это важно.} В задачах на совместную работу складываются не времена, а производительности. За один час вместе мастера выполняют сумму своих долей заказа. Затем эту общую производительность нужно умножить на два часа совместной работы, найти оставшуюся долю заказа и разделить её на производительность второго мастера.

\keyidea{Если весь заказ принят за единицу, то производительность мастера равна единице, делённой на время выполнения всего заказа. При совместной работе производительности складываются.}

\textbf{Исправление от последнего правильного шага.} Производительности мастеров:
\[
p_1=\frac{1}{5},\qquad p_2=\frac{1}{8}.
\]
Вместе за час они выполняют
\[
\frac{1}{5}+\frac{1}{8}=\frac{8+5}{40}=\frac{13}{40}
\]
заказа. За 2 часа выполнят
\[
2\cdot\frac{13}{40}=\frac{13}{20}.
\]
Останется
\[
1-\frac{13}{20}=\frac{7}{20}.
\]
Второй мастер выполняет по $1/8$ заказа в час, поэтому время на остаток:
\[
t=\frac{7/20}{1/8}=\frac{7}{20}\cdot8=\frac{14}{5}=2{,}8\ \text{ч}.
\]
Это 2 часа 48 минут.

\textbf{Полное правильное решение.} Примем весь заказ за 1. Первый мастер за один час делает $1/5$ заказа, второй --- $1/8$ заказа. Вместе за один час они выполняют
\[
\frac{1}{5}+\frac{1}{8}=\frac{13}{40}
\]
заказа. За два часа совместной работы выполнено
\[
2\cdot\frac{13}{40}=\frac{13}{20}.
\]
Оставшаяся часть равна
\[
1-\frac{13}{20}=\frac{7}{20}.
\]
После ухода первого мастера второй работает один со скоростью $1/8$ заказа в час. Поэтому ему требуется
\[
\frac{7}{20}:\frac{1}{8}=\frac{14}{5}\ \text{ч}=2{,}8\ \text{ч}.
\]
Десятые доли часа переводим в минуты:
\[
0{,}8\cdot60=48\ \text{мин}.
\]

\textbf{Проверка.} За 2 часа первый мастер выполняет $2/5$ заказа, второй --- $2/8=1/4$ заказа. Вместе это
\[
\frac{2}{5}+\frac{1}{4}=\frac{8}{20}+\frac{5}{20}=\frac{13}{20}.
\]
Остаётся $7/20$. За $14/5$ часа второй мастер выполнит
\[
\frac{1}{8}\cdot\frac{14}{5}=\frac{14}{40}=\frac{7}{20},
\]
то есть ровно остаток.

\finalresult{второму мастеру потребуется 2 часа 48 минут.}

\subsection*{Задача на закрепление}
Первый мастер выполняет заказ за 6 часов, второй --- за 9 часов. Они работают вместе 2 часа, после чего первый мастер уходит. Сколько времени потребуется второму мастеру, чтобы закончить оставшуюся работу?

\end{document}
